\documentclass[dvisvgm,tikz,border=5pt]{standalone}
\usepackage{amsmath,amssymb}
\usepackage{CJKutf8}
\usetikzlibrary{arrows.meta,positioning,calc,fit,decorations.pathreplacing}

\definecolor{themebg}{HTML}{2e362c}
\definecolor{themefg}{HTML}{edf4e8}
\definecolor{themegray}{HTML}{9ea897}
\definecolor{themecurve}{HTML}{7eb6ff}
\definecolor{themealert}{HTML}{ff7b7b}
\definecolor{themeamber}{HTML}{f5b942}

\begin{document}
\begin{CJK*}{UTF8}{gbsn}
\pagecolor{themebg}
\begin{tikzpicture}[font=\small,color=themefg,text=themefg,>=Stealth]
\tikzset{
  nodec/.style={circle,draw=themefg,fill=themebg,minimum size=7mm,inner sep=0pt,line width=.8pt},
  leaf/.style={rounded corners=2pt,draw=themefg,fill=themebg,minimum width=10mm,minimum height=7mm,inner sep=2pt},
  edge/.style={themegray,line width=.8pt},
  slot/.style={draw=themegray,minimum width=5.4mm,minimum height=5.5mm,inner sep=0pt,font=\scriptsize}
}

\node[anchor=west,font=\bfseries] at (0,6.7) {前缀码容量：短码会占掉它下面整棵后缀子树};
\node[anchor=west,font=\small,text=themegray] at (0,6.18)
  {二进制最大码长 $H=4$，已有码字 $1$、$01$；统一换算到深度 4 的 16 个叶槽。};

% trie sketch
\node[nodec] (root) at (3.0,5.0) {$\epsilon$};
\node[nodec] (z0) at (1.8,3.9) {0};
\node[leaf,draw=themealert,text=themealert] (z1) at (4.2,3.9) {1};
\draw[edge] (root)--node[above left,font=\scriptsize] {0}(z0);
\draw[edge] (root)--node[above right,font=\scriptsize] {1}(z1);
\node[nodec] (z00) at (.9,2.75) {00};
\node[leaf,draw=themeamber,text=themeamber] (z01) at (2.7,2.75) {01};
\draw[edge] (z0)--node[above left,font=\scriptsize] {0}(z00);
\draw[edge] (z0)--node[above right,font=\scriptsize] {1}(z01);
\node[font=\small,text=themealert,align=left,anchor=west] at (4.85,3.9)
  {$1$ 占 $2^{4-1}=8$ 个深度-4 槽};
\node[font=\small,text=themeamber,align=left,anchor=west] at (3.35,2.75)
  {$01$ 占 $2^{4-2}=4$ 个槽};
\node[font=\small,text=themecurve,align=left,anchor=west] at (1.45,1.8)
  {$00$ 子树还剩 4 个长度-4 码字位置};

% slot bar
\node[anchor=west,font=\bfseries] at (0,.95) {统一折算到深度 4：};
\foreach \i in {0,...,15}{
  \pgfmathsetmacro{\xx}{2.7+0.58*\i}
  \ifnum\i<8
    \node[slot,draw=themealert,text=themealert] at (\xx,.95) {\ifnum\i=3 $1$\fi};
  \else
    \ifnum\i<12
      \node[slot,draw=themeamber,text=themeamber] at (\xx,.95) {\ifnum\i=9 $01$\fi};
    \else
      \node[slot,draw=themecurve,text=themecurve] at (\xx,.95) {};
    \fi
  \fi
}
\draw[decorate,decoration={brace,mirror,amplitude=3pt},themealert,line width=.9pt]
  (2.42,.55) -- (6.98,.55) node[midway,below=5pt,font=\scriptsize,text=themealert] {$1$：8 槽};
\draw[decorate,decoration={brace,mirror,amplitude=3pt},themeamber,line width=.9pt]
  (7.05,.55) -- (9.37,.55) node[midway,below=5pt,font=\scriptsize,text=themeamber] {$01$：4 槽};
\draw[decorate,decoration={brace,mirror,amplitude=3pt},themecurve,line width=.9pt]
  (9.37,.55) -- (11.69,.55) node[midway,below=5pt,font=\scriptsize,text=themecurve] {剩余：4 槽};
\node[font=\small,text=themegray,anchor=west] at (2.65,-.15)
  {$16-8-4=4$：最多再增加 4 个长度为 4 的码字；若问总字符数，则为 $2+4=6$。};

\end{tikzpicture}
\end{CJK*}
\end{document}
